\section{Discussion}
\label{sec:discussion}
Pakistan-specific activity and process boundaries determine the inventory's contribution. Crop, nutrient and census-anchored herd controls support a national reconstruction whose source pathways and geographic support remain inspectable. The common nitrogen ledger permits differences in NH$_3$, NO$_x$ and N$_2$O to be explained by pathway or factor rather than hidden in incompatible accounts.

The retained dry-climate soil factor is consequential and conditional. Irrigated cropland, flooded rice and grazing have different factor applicability, so a national climate label does not resolve every pathway. Aggregate and climate-factor sensitivities show the consequence without fitting to FAOSTAT or national reporting. The same distinction applies to manure-system shares: material closure checks the calculation but does not establish the chosen management distribution.

Independent district crop records evaluate selected spatial patterns. Inventories constructed from the same livestock counts or cropland surfaces do not independently confirm those allocations. Coarse support downscaled to 0.05$^\circ$ retains proxy limitations. Source-matched public HTAP comparisons supplement national reporting and global inventories, with provider and shared-input provenance explicit.

The common core's latest annual year and its detailed reference year have different roles. FY2024--25 uses supportable annual activity; it does not acquire new monthly application dates, herd enumeration or fine-scale source geography automatically. Fitness for use therefore depends on species, emitting process, year and configuration as well as total mass. Atmospheric concentrations, exposure and health effects require additional evidence and are outside this inventory.
